\documentclass[11pt]{article}
\usepackage[margin=1in]{geometry}
\usepackage{amsmath,amssymb}
\usepackage{booktabs}
\usepackage{graphicx}
\usepackage{hyperref}
\usepackage{xcolor}
\usepackage{natbib}
\usepackage{enumitem}
\setlist{nosep}

\newcommand{\eps}{\varepsilon}
\newcommand{\epspert}{\eps_\mathrm{pert}}
\newcommand{\rhogeo}{\rho_\mathrm{geo}}
\newcommand{\rhohol}{\rho_\mathrm{hol}}
\newcommand{\Gr}{\mathrm{Gr}}

\title{Separating Signal from Noise in Geometric\\
  Cross-Sensor Metrics: A Pre-Registered Case Study}
\author{Elliot Tower}
\date{July 2026 --- Draft v1}

\begin{document}
\maketitle

% ================================================================
\begin{abstract}
Grassmannian geodesic distance between PCA subspaces offers a
parameter-free measure of cross-sensor consistency: when multiple
sensors observe the same latent process, their estimated subspaces
should align, and divergence signals genuine disagreement about the
underlying signal.  Validating this metric, however, requires care.
We present a case study in which a naive single-parameter simulation
design confounded observation noise with embedding perturbation,
producing spurious positive results that survived standard
significance testing.  Through four pre-registered iterations---each
corrected after negative controls exposed a specific flaw---we
developed a two-axis experimental design that cleanly separates
noise from perturbation.  The resulting degradation curve
characterizes the operating regime: at low noise
($\sigma = 1.0$), geodesic distance reliably detects cross-sensor
perturbation (median $\rhogeo = 0.52$, 95\% CI excludes zero,
62\% per-seed detection); at high noise ($\sigma = 4.5$), the
signal is fully attenuated (median $\rhogeo = 0.15$, 0\%
per-seed detection).  We also report a negative result for loop
holonomy, a candidate severity metric: it is genuinely
non-monotone in perturbation magnitude, functioning as a
saturating onset detector rather than a graded tracker.  The
two-axis design principle and degradation-curve methodology
transfer to any setting where geometric cross-sensor metrics
require validation.
\end{abstract}

% ================================================================
\section{Introduction}
\label{sec:intro}

When multiple sensors observe the same latent process, a natural
question arises: do their measurements agree?  Agreement at the
level of summary statistics (means, variances) can mask structural
disagreement about which directions in feature space carry signal.
Grassmannian geodesic distance---the arc length between subspaces
estimated by PCA---captures this directional information without
learned parameters, offering a geometric measure of cross-sensor
consistency
\textcolor{red}{[CITE: Edelman et al.\ 1998, geometry of
Grassmannian]}.

The appeal is clear: a single number quantifying how much two
sensors' measurement bases have diverged, computable from data
alone.  The difficulty is validation.  On real multi-sensor data,
ground truth is unavailable---we cannot directly observe whether a
large geodesic distance reflects genuine inconsistency or
measurement noise.  Simulation studies fill this gap, but they
introduce a subtler problem: the simulation design itself can
confound the effects it aims to isolate.

This paper documents such a confound and the iterative process that
uncovered it.  We began with a single-parameter design in which one
variable $\eps$ simultaneously controlled observation noise and
cross-sensor embedding perturbation.  Under this design, geodesic
distance increased monotonically with $\eps$---an apparently clean
signal.  Negative controls revealed the truth: the monotonicity was
driven primarily by noise scaling, not perturbation.  Identical
sensors (no embedding divergence) produced \emph{stronger}
monotonicity than genuinely different sensors, because the noise
ramp dominated the geodesic metric.

Correcting this required separating the two axes: observation noise
$\sigma$ (held constant within each experimental sweep) and
embedding perturbation $\epspert$ (varied independently).  The
separation was not obvious from the initial design.  It emerged
through four pre-registered iterations (Section~\ref{sec:confound}),
each frozen under a commit SHA before any confirmatory data was
generated.

The resulting experiment (Section~\ref{sec:results}) yields three
contributions:

\begin{enumerate}
  \item A \textbf{two-axis design principle} for validating
    geometric cross-sensor metrics, applicable beyond the specific
    simulation studied here.
  \item A \textbf{degradation curve} characterizing the
    operating-regime boundary: the noise level at which geodesic
    distance ceases to reliably detect perturbation.
  \item A \textbf{negative result} for loop holonomy: the composed
    parallel transport around a sensor loop is non-monotone in
    perturbation, confirmed by power analysis as a property of the
    metric (not insufficient sample size).
\end{enumerate}

The paper's broader point is methodological.  Geometric metrics are
attractive precisely because they are parameter-free and
deterministic.  These properties do not exempt them from validation,
and validation requires experimental designs that are themselves
free of confounds.  Pre-registered negative controls are an
effective tool for catching such confounds, even in simulation
settings where the data-generating process is fully known.

% ================================================================
\section{The confound discovery}
\label{sec:confound}

This section traces the four design iterations that preceded the
confirmatory experiment.  Each iteration discovered a flaw through
pre-registered controls, motivating a redesign.  We present them
chronologically because the sequence itself illustrates how
confounds hide in apparently reasonable designs.

\subsection{Version 1--2: The wrong null and the wrong signal}
\label{sec:v1v2}

The initial design simulated three sensors in $\mathbb{R}^{100}$
observing a shared $k$-dimensional latent signal ($k = 3$) through
sensor-specific linear embeddings.  A single parameter $\eps \in
[0, 1]$ controlled the ``decoherence level,'' simultaneously
scaling observation noise ($\sigma = 1.0 + 3.5\eps$) and
sensor-specific embedding perturbation ($U_s(\eps) =
\mathrm{QR}(U_\mathrm{base} + \eps \cdot 1.5 \cdot P_s)$, where
$P_s$ is a random perturbation matrix unique to each sensor).

Two problems emerged before any confirmatory run:

\textbf{Broken permutation null.}  The initial null distribution
permuted rows within each sensor, then recomputed PCA subspaces.
Row permutation does not change the covariance matrix, so PCA
returns identical subspaces regardless of the permutation.  The
null distribution collapsed to a single point (empirical standard
deviation: $0.000$), making the test trivially significant.  A
pool-and-resplit null replaced it.

\textbf{Noise heterogeneity masquerading as inconsistency.}  The
three sensors differed only in additive noise scaling, not in
embedding direction.  The geodesic metric fired on distributional
differences (different variance patterns), not on genuine geometric
disagreement about the latent signal.  A redesign introduced
sensor-specific embedding perturbations so that at $\eps = 0$ all
sensors share the same measurement basis, and at $\eps > 0$ each
sensor's basis tilts in a different direction.

\subsection{Version 3: The single-parameter confound}
\label{sec:v3}

The redesigned data model (version~3) appeared sound: genuinely
different embeddings, a functional null, controlled signal
strength.  Five seeds confirmed the geodesic metric's monotonicity
(Spearman $\rho > 0.4$ in 4 of 5 seeds).  A deeper analysis,
however, revealed two problems.

\textbf{Loop holonomy is non-monotone.}  The composed parallel
transport around the three-sensor loop---a candidate severity
metric---showed Spearman $\rho$ between $0.01$ and $0.40$ across
seeds (none significant at $p < 0.05$).  Pairwise edge geodesics
were monotone ($\rho \approx 0.5$), but the loop composition
introduced cancellations that destroyed the signal.  A power
analysis over sample sizes $N \in \{200, 500, 1000, 2000\}$
confirmed this is a property of the metric, not a sample-size
limitation: Spearman $\rho$ plateaued well below $+1$ and
confidence bands did not narrow.

\textbf{The single-parameter confound.}  The definitive failure
came from a negative control.  Three sensors with \emph{identical}
embedding perturbation ($P_0 = P_1 = P_2$, no geometric
inconsistency) should produce near-zero differential signal
($\Delta\rho = \rho_\mathrm{edge} - \rho_\mathrm{loop} \approx
0$).  Instead, the negative control produced $\Delta\rho \approx
+0.60$---\emph{larger} than the $\Delta\rho \approx +0.35$
observed under genuine inconsistency.

The mechanism: when all sensors share the same embedding, the
geodesic metric tracks observation noise scaling ($\sigma = 1.0 +
3.5\eps$) with near-perfect monotonicity ($\rho \approx 0.96$),
because increasing noise degrades PCA subspace estimation
symmetrically across sensors.  Genuine embedding perturbation
\emph{disrupts} this clean noise-driven monotonicity, producing a
lower (not higher) correlation.  The single-parameter design made
every downstream metric a proxy for noise scaling, not cross-sensor
inconsistency.

\subsection{Version 4: The two-axis design}
\label{sec:v4}

The fix separates the single $\eps$ into two orthogonal axes:

\begin{itemize}
  \item \textbf{Embedding perturbation} $\epspert \in [0, 1]$
    controls cross-sensor divergence only: $U_s(\epspert) =
    \mathrm{QR}(U_\mathrm{base} + \epspert \cdot 1.5 \cdot P_s)$.
  \item \textbf{Observation noise} $\sigma$ is held constant within
    each experimental sweep.
\end{itemize}

Independence of the generative axes was verified: sensor embeddings
depend only on $\epspert$, observation noise depends only on
$\sigma$, and no parameter is shared.  A known measurement-level
interaction remains: noise degrades PCA subspace estimation, which
attenuates the perturbation signal at high $\sigma$.  This is
expected physics (noise masks signal), not a generative-model
confound.  The noise grid characterizes this attenuation rather
than hiding it.

The primary endpoint sweeps $\epspert$ at each of four fixed noise
levels ($\sigma \in \{1.0, 2.0, 3.0, 4.5\}$), with the
confirmatory test at $\sigma = 1.0$ (pre-committed as the most
favorable operating point).  Separate seed pools (50 primary, 50
negative control, 50 noise-axis control) ensure no data reuse.

% ================================================================
\section{Methods}
\label{sec:methods}

\subsection{Data model}

Three sensors observe a shared latent process in
$\mathbb{R}^{100}$.  The latent data comprise $N = 200$ samples
per class (two classes) drawn from $\mathbb{R}^3$ with class
signal scaled by $\lambda = 5.0$.  Each sensor projects through
$U_s(\epspert) = \mathrm{QR}(U_\mathrm{base} + \epspert \cdot 1.5
\cdot P_s)$, where $U_\mathrm{base} \in \mathbb{R}^{100 \times 3}$
is a shared base embedding and $P_s \in \mathbb{R}^{100 \times 3}$
is a sensor-specific perturbation matrix.  Observations are
$X_s = Z \cdot U_s^\top + \mathcal{N}(0, \sigma^2 I)$.

At $\epspert = 0$, all sensors share $U_\mathrm{base}$: no
cross-sensor inconsistency exists.  At $\epspert > 0$, each
sensor's measurement basis tilts in a different direction,
creating genuine geometric inconsistency detectable on the
Grassmannian.

\subsection{Grassmannian geodesic distance}

For each sensor, PCA extracts the top-$k$ subspace ($k = 3$) from
the observation matrix.  Two subspaces $\mathcal{U},
\mathcal{V} \in \Gr(k, d)$ are compared via the geodesic distance:
\[
  d_{\Gr}(\mathcal{U}, \mathcal{V}) =
    \|\boldsymbol{\theta}\|_2
    = \sqrt{\theta_1^2 + \cdots + \theta_k^2},
\]
where $\theta_1, \ldots, \theta_k$ are the principal angles between
$\mathcal{U}$ and $\mathcal{V}$
\textcolor{red}{[CITE: Ye \& Lim 2016, Grassmannian distances]}.
The mean of all three pairwise geodesic distances serves as the
per-step summary statistic.

\subsection{Loop holonomy}

Parallel transport around the sensor loop $0 \to 1 \to 2 \to 0$
is computed as $T = U_0^\top U_1 \cdot U_1^\top U_2 \cdot
U_2^\top U_0$, where $U_s$ are the $d \times k$ PCA basis matrices.
The holonomy is $\|I_k - T\|_F$, measuring the failure of a vector
to return to itself after transport around the loop.

\subsection{Statistical framework}

\textbf{Primary statistic.}  For each seed at each fixed noise
level, compute $\rhogeo = \mathrm{Spearman}(\text{mean geodesic},
\, \epspert)$ over the 20-step sweep $\epspert \in [0, 1]$.

\textbf{Aggregate inference.}  The median $\rhogeo$ across 50 seeds
and its 95\% bootstrap confidence interval (10{,}000 resamples,
bootstrap RNG seed = 42) constitute the primary summary.  The
confirmatory test at $\sigma = 1.0$ is confirmed if the CI excludes
zero ($\alpha = 0.05$, single test, no multiplicity correction).

\textbf{Per-seed inference.}  Each seed's $p$-value from the
Spearman test is corrected within each noise level using
Benjamini--Hochberg FDR at $\alpha = 0.05$.  The fraction of
significant seeds characterizes individual-seed reliability.

\textbf{Negative control.}  Three sensors with identical
perturbation ($P_0 = P_1 = P_2$) at each noise level; the
bootstrap CI must include zero.

\textbf{Noise-axis control.}  Sweep $\sigma$ from 1.0 to 4.5 (20
steps) at $\epspert = 0$ (no perturbation).  The strongly positive
$\rhogeo$ expected here documents the nuisance axis rather than
hiding it.

\textbf{Degradation curve.}  Spearman correlation of median
$\rhogeo$ across the four noise levels.

\subsection{Pre-registration protocol}

All parameters, seed assignments, and statistical procedures were
frozen under commit SHA \texttt{a792e25} before any confirmatory
seeds were generated.  Sanity seeds (10, used for parameter
calibration) were excluded from the confirmatory analysis.
Primary, negative control, and noise-axis control each used 50
non-overlapping seeds.

% ================================================================
\section{Results}
\label{sec:results}

\subsection{Confirmatory test ($\sigma = 1.0$)}

At the pre-committed most-favorable noise level, the median
$\rhogeo$ across 50 primary seeds was $0.518$ (95\% CI: $[0.462,
0.576]$), excluding zero.  The negative control median was $0.050$
($[-0.065, 0.137]$), including zero as expected
(Figure~\ref{fig:degradation}a).

\subsection{Degradation curve}

Detection strength decreased monotonically with noise level
(Table~\ref{tab:degradation}, Figure~\ref{fig:degradation}a).  The
aggregate CI excluded zero at all four levels, but per-seed
detection power collapsed rapidly: 62\% at $\sigma = 1.0$, 20\%
at $\sigma = 2.0$, 2\% at $\sigma = 3.0$, and 0\% at
$\sigma = 4.5$ (Figure~\ref{fig:degradation}b).

\begin{table}[h]
\centering
\caption{Degradation of geodesic detection across noise levels.
  Aggregate CI excludes zero at all levels, but per-seed power
  collapses above $\sigma \approx 2$.}
\label{tab:degradation}
\begin{tabular}{@{}ccccc@{}}
\toprule
$\sigma$ & Median $\rhogeo$ & 95\% CI & Excl.\ 0? & BH sig. \\
\midrule
1.0 & 0.518 & $[0.462, 0.576]$ & Yes & 31/50 \\
2.0 & 0.438 & $[0.389, 0.495]$ & Yes & 10/50 \\
3.0 & 0.325 & $[0.284, 0.410]$ & Yes & \phantom{0}1/50 \\
4.5 & 0.153 & $[0.077, 0.198]$ & Yes & \phantom{0}0/50 \\
\bottomrule
\end{tabular}
\end{table}

The distinction between aggregate and per-seed significance is
important.  Pooling 50 seeds provides sufficient power to detect a
small residual signal even at $\sigma = 4.5$, but no individual
seed at that noise level achieves significance after BH correction.
For practical deployment, the per-seed rate is the relevant measure:
a single dataset from a single experiment must stand on its own.
By this criterion, the practical operating boundary lies near
$\sigma \approx 2$, where per-seed power drops below 50\%.

The Spearman correlation across the four degradation-curve points is
$-1.0$, confirming strict monotonic decay.  We note that four
strictly decreasing points guarantee $\rho = -1.0$; a finer noise
grid would better characterize the functional form.

\subsection{Negative controls}

At all four noise levels, the negative control (identical
perturbation) produced median $\rhogeo$ with bootstrap CIs
including zero: $0.050$ at $\sigma = 1.0$, $-0.022$ at
$\sigma = 2.0$, $-0.004$ at $\sigma = 3.0$, and $-0.010$ at
$\sigma = 4.5$.  The two-axis design eliminates the spurious
signal that contaminated the single-parameter design.

\subsection{Noise-axis control}

Sweeping $\sigma$ at $\epspert = 0$ yielded median $\rhogeo =
0.966$ ($[0.958, 0.973]$), confirming that geodesic distance is
highly sensitive to observation noise when noise varies.  This
characterizes the nuisance axis: the very effect that confounded
the single-parameter design (Section~\ref{sec:v3}) is documented
and accounted for by holding $\sigma$ fixed in the primary
endpoint.

\subsection{Loop holonomy (exploratory)}

At $\sigma = 1.0$, loop holonomy showed median $\rhohol = 0.263$
($[0.157, 0.352]$)---positive but substantially weaker than the
geodesic signal ($\rhogeo = 0.518$).  At $\sigma = 4.5$,
$\rhohol = -0.039$, consistent with zero.  Combined with the
version~3 finding that holonomy is non-monotone in perturbation
even under fixed noise (Section~\ref{sec:v3}), these results
confirm that loop holonomy functions as a saturating onset
detector, not a graded severity tracker.  Pairwise geodesic
distance is the appropriate metric for quantitative assessment of
cross-sensor divergence; loop holonomy may retain value as a
qualitative indicator of \emph{whether} inconsistency exists, but
not \emph{how much}.

% ================================================================
\section{Discussion}
\label{sec:discussion}

\subsection{Transferable design principles}

The two-axis separation---holding one effect constant while varying
another---is a general principle for validating any metric that
responds to multiple sources of variation.  Geodesic distance
responds to both observation noise (through PCA subspace
estimation error) and embedding perturbation (through genuine
subspace divergence).  A single-parameter design that varies both
simultaneously cannot attribute metric changes to either cause.
This confound is not exotic: any simulation in which ``noise level''
and ``effect size'' covary through a shared parameter is
susceptible.

The negative control design transfers directly.  Running the full
experimental protocol on data where the target effect is absent
(identical perturbation) is the most reliable way to detect
confounds, because it tests the entire analysis pipeline---not
just the metric's mathematical properties in isolation.

\subsection{The degradation curve as operating-regime characterization}

Rather than reporting a single pass/fail verdict, the degradation
curve maps the noise-detection tradeoff continuously.  The
aggregate signal persists to $\sigma = 4.5$, but per-seed power
collapses above $\sigma \approx 2$.  This distinction matters for
practitioners: an aggregate analysis pooling many datasets may
detect subtle inconsistency that no individual dataset can support.
Reporting both aggregate and per-seed power avoids over-selling
the metric's reliability.

The degradation curve format---median effect size with bootstrap CI
at each noise level, alongside per-seed detection power---is
reusable for any metric validation study.  It separates the
question ``does the metric detect the effect in principle?'' from
``does it detect the effect in a single experiment?''

\subsection{Holonomy: a negative result}

Loop holonomy around a multi-sensor loop was initially proposed as
a severity metric that could detect inconsistency at lower
thresholds than pairwise geodesic distance.  The evidence refutes
this: holonomy is non-monotone in perturbation, confirmed across
multiple sample sizes and seed pools.  The composed parallel
transport introduces cancellations when sensors perturb in
different directions, destroying the graded signal present in
individual edges.

This negative result has practical value.  Loop holonomy appears in
proposals for multi-sensor consistency checking
\textcolor{red}{[CITE: relevant holonomy papers if any]}, and the
non-monotonicity reported here constrains its appropriate use to
binary detection (consistent vs.\ inconsistent), not quantitative
severity assessment.

\subsection{Limitations}

The data model is linear (latent-to-observation projection),
low-dimensional ($d = 100$, $k = 3$), and Gaussian-noised.  Real
multi-sensor systems involve nonlinear transductions, higher
ambient dimensions, and non-Gaussian noise distributions.  The
specific numbers in the degradation curve (e.g., the practical
boundary at $\sigma \approx 2$) depend on these choices and should
not be taken as universal constants.

What transfers is the methodology: the two-axis design principle,
the negative-control strategy, and the degradation-curve format.
These are independent of the simulation's specifics.

The degradation curve uses four noise levels.  A finer grid (8--12
levels) would better characterize the functional form of the
decay and locate the per-seed power boundary more precisely.

\subsection{The value of iterative pre-registration}

The confounds documented in Section~\ref{sec:confound} were not
discovered by mathematical analysis of the data model.  They were
discovered by running the experiment under pre-registered controls
and observing unexpected results.  The permutation null's collapse,
the noise-heterogeneity artifact, and the single-parameter
confound were all invisible until the data contradicted the
predictions.

Pre-registration in simulation settings may seem unnecessary---the
data-generating process is fully known, so there is nothing to
``pre-register against.''  Our experience suggests otherwise.
Knowing the generative model does not guarantee understanding its
interaction with the analysis pipeline.  The single-parameter
confound arose not from the generators but from the coupling
between noise scaling and PCA subspace estimation error, an
interaction that was not anticipated from the model specification.

Freezing parameters, seeds, and analysis plans under a commit SHA
before each confirmatory run imposed a discipline that repeatedly
surfaced these problems.  Without it, the version~3 result---a
clean-looking positive signal driven entirely by noise
scaling---would have been reported as a confirmation.

% ================================================================
\section*{Figures}

\textbf{Figure~1} (Fig.~\ref{fig:degradation}): Degradation curve.
(a)~Median $\rhogeo$ vs.\ $\sigma$ with bootstrap CIs, holonomy
overlay, negative control.  (b)~Per-seed BH detection rate.
[\texttt{b2\_degradation\_curve.pdf}]

\textbf{Figure~2}: Seed trajectories at $\sigma = 1.0$.
(a)~Raw geodesic vs.\ $\epspert$ for 5 primary and 5 negative
control seeds.  (b)~Distribution of $\rhogeo$ across all 50 seeds.
[\texttt{b2\_seed\_trajectories.pdf}]

\textbf{Figure~3}: Experimental design comparison.
(a)~Confounded single-$\eps$ design: one parameter drives both
noise and perturbation into the geodesic.  (b)~Orthogonal two-axis
design: $\sigma$ held constant, $\epspert$ varied independently.
[\texttt{b2\_confound\_schematic.pdf}]

\textbf{Figure~4}: Noise-axis control.
(a)~Raw geodesic vs.\ $\sigma$ for 5 seeds at $\epspert = 0$.
(b)~Distribution of $\rhogeo$ across all 50 seeds (median $= 0.966$).
[\texttt{b2\_noise\_axis\_control.pdf}]

\label{fig:degradation}

% ================================================================
% \bibliographystyle{plainnat}
% \bibliography{refs}

\end{document}
